\chapter{Programmation en langage C}

\textit{Le langage C est un langage de programmation impératif, structuré et largement utilisé pour la programmation système et embarquée. Ce chapitre présente les bases essentielles pour écrire, comprendre et corriger des programmes en C.}

\section{L'essentiel du cours}

\subsection{Structure d'un programme C}

Un programme C est composé de fonctions. La fonction \texttt{main()} est le point d'entrée du programme. Les instructions se terminent par un point-virgule.

\begin{ccode}
#include <stdio.h>   // Inclusion de la bibliothèque standard d'entrée/sortie

int main() {
    printf("Bonjour le Cameroun !\n");
    return 0;
}
\end{ccode}

\begin{aretenir}
Tout programme C commence par l'exécution de la fonction \texttt{main()}. Le \texttt{return 0;} indique que le programme s'est terminé avec succès.
\end{aretenir}

\subsection{Types et variables}

Une variable est un espace mémoire nommé qui stocke une valeur d'un type donné.

\begin{definition}[Types de base en C]
\begin{itemize}
    \item \texttt{int} : entier (ex: 25, -3, 0)
    \item \texttt{float} : nombre à virgule flottante (ex: 3.14, -0.5)
    \item \texttt{double} : nombre à virgule flottante double précision
    \item \texttt{char} : caractère (ex: 'A', '7', '\textbackslash n')
\end{itemize}
\end{definition}

\begin{ccode}
int age = 17;
float note = 14.5;
char lettre = 'B';
\end{ccode}

\subsection{Entrées/Sorties (printf, scanf)}

\begin{itemize}
    \item \texttt{printf("format", variables)} : affiche du texte à l'écran.
    \item \texttt{scanf("format", \&variables)} : lit une valeur saisie au clavier.
\end{itemize}

\begin{ccode}
int nombre;
printf("Entrez un nombre : ");
scanf("%d", &nombre);
printf("Vous avez saisi : %d\n", nombre);
\end{ccode}

\begin{aretenir}
Les formats courants : \texttt{\%d} pour \texttt{int}, \texttt{\%f} pour \texttt{float}, \texttt{\%c} pour \texttt{char}. N'oubliez pas le \texttt{\&} devant la variable dans \texttt{scanf}.
\end{aretenir}

\subsection{Opérateurs}

\begin{itemize}
    \item \textbf{Arithmétiques} : \texttt{+}, \texttt{-}, \texttt{*}, \texttt{/}, \texttt{\%} (modulo)
    \item \textbf{Comparaison} : \texttt{==}, \texttt{!=}, \texttt{<}, \texttt{>}, \texttt{<=}, \texttt{>=}
    \item \textbf{Logiques} : \texttt{\&\&} (ET), \texttt{||} (OU), \texttt{!} (NON)
    \item \textbf{Incrémentation} : \texttt{++} (ajoute 1), \texttt{--} (enlève 1)
\end{itemize}

\subsection{Structures de contrôle}

\subsubsection{Conditionnelles : if/else}

\begin{ccode}
if (note >= 10) {
    printf("Admis\n");
} else {
    printf("Échec\n");
}
\end{ccode}

\subsubsection{Boucles : for et while}

\begin{ccode}
// Boucle for : pour i de 1 à 5
for (int i = 1; i <= 5; i++) {
    printf("%d ", i);
}

// Boucle while
int j = 1;
while (j <= 5) {
    printf("%d ", j);
    j++;
}
\end{ccode}

\subsection{Tableaux}

Un tableau stocke plusieurs valeurs du même type.

\begin{ccode}
int notes[5] = {12, 15, 8, 10, 18};
notes[2] = 14;  // Modification de la 3e case (indice 2)
printf("Note 1 : %d\n", notes[0]);
\end{ccode}

\begin{aretenir}
L'indice du premier élément est 0. Un tableau de taille \texttt{n} a des indices de 0 à \texttt{n-1}.
\end{aretenir}

\subsection{Fonctions}

Une fonction est un bloc de code réutilisable.

\begin{ccode}
// Définition d'une fonction
int carre(int x) {
    return x * x;
}

int main() {
    int resultat = carre(5);
    printf("5 au carré = %d\n", resultat);
    return 0;
}
\end{ccode}

\begin{propriete}
Une fonction peut avoir des paramètres (entrées) et retourner une valeur avec \texttt{return}. Si elle ne retourne rien, on utilise \texttt{void}.
\end{propriete}

\section{Exercices}

\rubrique{Structure et affichage}

\exo{}\diff{1}
Écrire un programme C qui affiche "Bienvenue en Terminale C" à l'écran.

\exo{}\diff{1}
Écrire un programme C qui déclare une variable \texttt{annee} de type \texttt{int} initialisée à 2025, puis l'affiche.

\exo{}\diff{1}
Compléter le programme suivant pour qu'il affiche la somme de deux entiers saisis par l'utilisateur.

\begin{ccode}
#include <stdio.h>
int main() {
    int a, b, somme;
    printf("Entrez deux nombres : ");
    // À compléter
    return 0;
}
\end{ccode}

\rubrique{Opérateurs et calculs}

\exo{}\diff{1}
Écrire un programme qui lit un nombre entier et affiche son carré et son cube.

\exo{}\diff{1}
Écrire un programme qui lit deux entiers et affiche leur quotient et leur reste (division euclidienne).

\exo{}\diff{2}
Écrire un programme qui lit un nombre entier à trois chiffres et affiche la somme de ses chiffres. (Exemple : 345 → 3+4+5 = 12)

\rubrique{Structures conditionnelles}

\exo{}\diff{1}
Écrire un programme qui lit un entier et affiche "Pair" s'il est pair, "Impair" sinon.

\exo{}\diff{1}
Écrire un programme qui lit trois notes (sur 20) et affiche la moyenne. Si la moyenne est supérieure ou égale à 10, afficher "Admis", sinon "Échec".

\exo{}\diff{2}
Écrire un programme qui lit un caractère et détermine s'il s'agit d'une lettre majuscule, d'une lettre minuscule ou d'un chiffre.

\exo{}\sujetbac[2023, C]{}\diff{2}
Un élève souhaite écrire un programme qui lit l'âge d'une personne et affiche "Mineur" si l'âge est inférieur à 18, "Majeur" sinon. Corriger les erreurs dans le code suivant :

\begin{ccode}
#include <stdio.h>
int main() {
    int age;
    printf("Entrez votre age : ");
    scanf("%d", age);
    if (age < 18) {
        printf("Mineur\n");
    else {
        printf("Majeur\n");
    }
    return 0;
}
\end{ccode}

\rubrique{Boucles}

\exo{}\diff{1}
Écrire un programme qui affiche les nombres de 1 à 10 à l'aide d'une boucle \texttt{for}.

\exo{}\diff{1}
Écrire un programme qui affiche les nombres pairs de 2 à 20.

\exo{}\diff{2}
Écrire un programme qui lit un entier \texttt{n} et calcule la somme des \texttt{n} premiers entiers naturels (1+2+...+n).

\exo{}\diff{2}
Écrire un programme qui lit un entier \texttt{n} et affiche sa table de multiplication (de 1 à 10).

\exo{}\sujetbac[2022, D]{}\diff{2}
On donne le programme suivant. Quelle sera la sortie ?

\begin{ccode}
#include <stdio.h>
int main() {
    int i, s = 0;
    for (i = 1; i <= 5; i++) {
        s = s + i;
    }
    printf("s = %d\n", s);
    return 0;
}
\end{ccode}

\exo{}\diff{3}
Écrire un programme qui lit un entier \texttt{n} et affiche s'il est premier ou non.

\rubrique{Tableaux}

\exo{}\diff{1}
Écrire un programme qui déclare un tableau de 5 entiers, les initialise avec les valeurs 10, 20, 30, 40, 50, puis les affiche.

\exo{}\diff{2}
Écrire un programme qui lit 10 notes (sur 20) dans un tableau, puis calcule et affiche la moyenne.

\exo{}\diff{2}
Écrire un programme qui lit 5 entiers dans un tableau, puis affiche le plus grand élément et sa position.

\exo{}\sujetbac[2021, C]{}\diff{3}
On considère le programme suivant. Compléter les pointillés pour qu'il affiche le nombre d'occurrences d'une valeur \texttt{v} lue dans un tableau de 10 entiers.

\begin{ccode}
#include <stdio.h>
int main() {
    int tab[10], v, i, compteur = 0;
    printf("Entrez 10 entiers : ");
    for (i = 0; i < 10; i++) {
        scanf("%d", &tab[i]);
    }
    printf("Entrez la valeur à chercher : ");
    scanf("%d", &v);
    for (i = 0; i < 10; i++) {
        // À compléter
    }
    printf("La valeur %d apparaît %d fois.\n", v, compteur);
    return 0;
}
\end{ccode}

\rubrique{Fonctions}

\exo{}\diff{1}
Écrire une fonction \texttt{max2} qui prend deux entiers et retourne le plus grand.

\exo{}\diff{2}
Écrire une fonction \texttt{estPair} qui prend un entier et retourne 1 s'il est pair, 0 sinon. L'utiliser dans un programme principal.

\exo{}\diff{2}
Écrire une fonction \texttt{sommeTableau} qui prend un tableau d'entiers et sa taille, et retourne la somme de ses éléments.

\exo{}\sujetbac[2020, E]{}\diff{3}
On veut écrire une fonction \texttt{factorial} qui calcule la factorielle d'un entier \texttt{n} (n! = 1×2×...×n). Proposer le code de cette fonction, puis l'utiliser dans un programme qui lit un entier et affiche sa factorielle.

\rubrique{Programmes complets}

\exo{}\diff{2}
Écrire un programme complet qui lit un entier \texttt{n} et affiche tous ses diviseurs.

\exo{}\diff{3}
Écrire un programme qui lit 10 entiers dans un tableau, puis les trie par ordre croissant (tri à bulles).

\exo{}\sujetbac[2019, C]{}\diff{3}
Un magasin de Yaoundé applique une remise de 10\% sur le montant total des achats si celui-ci dépasse 50 000 FCFA. Écrire un programme qui lit le montant des achats et affiche le montant à payer après remise éventuelle.

\section{Corrigés}

\corrige{de l'exercice 1}
\begin{ccode}
#include <stdio.h>
int main() {
    printf("Bienvenue en Terminale C\n");
    return 0;
}
\end{ccode}

\corrige{de l'exercice 2}
\begin{ccode}
#include <stdio.h>
int main() {
    int annee = 2025;
    printf("Année : %d\n", annee);
    return 0;
}
\end{ccode}

\corrige{de l'exercice 3}
\begin{ccode}
#include <stdio.h>
int main() {
    int a, b, somme;
    printf("Entrez deux nombres : ");
    scanf("%d %d", &a, &b);
    somme = a + b;
    printf("Somme = %d\n", somme);
    return 0;
}
\end{ccode}

\corrige{de l'exercice 4}
\begin{ccode}
#include <stdio.h>
int main() {
    int n;
    printf("Entrez un entier : ");
    scanf("%d", &n);
    printf("Carré = %d\n", n * n);
    printf("Cube = %d\n", n * n * n);
    return 0;
}
\end{ccode}

\corrige{de l'exercice 5}
\begin{ccode}
#include <stdio.h>
int main() {
    int a, b;
    printf("Entrez deux entiers : ");
    scanf("%d %d", &a, &b);
    if (b != 0) {
        printf("Quotient = %d\n", a / b);
        printf("Reste = %d\n", a % b);
    } else {
        printf("Division par zéro impossible.\n");
    }
    return 0;
}
\end{ccode}

\corrige{de l'exercice 6}
\begin{ccode}
#include <stdio.h>
int main() {
    int n, centaines, dizaines, unites, somme;
    printf("Entrez un entier à trois chiffres : ");
    scanf("%d", &n);
    centaines = n / 100;
    dizaines = (n / 10) % 10;
    unites = n % 10;
    somme = centaines + dizaines + unites;
    printf("Somme des chiffres = %d\n", somme);
    return 0;
}
\end{ccode}

\corrige{de l'exercice 7}
\begin{ccode}
#include <stdio.h>
int main() {
    int n;
    printf("Entrez un entier : ");
    scanf("%d", &n);
    if (n % 2 == 0) {
        printf("Pair\n");
    } else {
        printf("Impair\n");
    }
    return 0;
}
\end{ccode}

\corrige{de l'exercice 8}
\begin{ccode}
#include <stdio.h>
int main() {
    float n1, n2, n3, moyenne;
    printf("Entrez trois notes : ");
    scanf("%f %f %f", &n1, &n2, &n3);
    moyenne = (n1 + n2 + n3) / 3;
    printf("Moyenne = %.2f\n", moyenne);
    if (moyenne >= 10) {
        printf("Admis\n");
    } else {
        printf("Échec\n");
    }
    return 0;
}
\end{ccode}

\corrige{de l'exercice 9}
\begin{ccode}
#include <stdio.h>
int main() {
    char c;
    printf("Entrez un caractère : ");
    scanf(" %c", &c);
    if (c >= 'A' && c <= 'Z') {
        printf("Majuscule\n");
    } else if (c >= 'a' && c <= 'z') {
        printf("Minuscule\n");
    } else if (c >= '0' && c <= '9') {
        printf("Chiffre\n");
    } else {
        printf("Autre caractère\n");
    }
    return 0;
}
\end{ccode}

\corrige{de l'exercice 10}
Erreurs corrigées :
\begin{itemize}
    \item Dans \texttt{scanf}, il manquait le \texttt{\&} devant \texttt{age}.
    \item Il manquait une accolade fermante \texttt{\}} avant le \texttt{else}.
\end{itemize}
Code corrigé :
\begin{ccode}
#include <stdio.h>
int main() {
    int age;
    printf("Entrez votre age : ");
    scanf("%d", &age);
    if (age < 18) {
        printf("Mineur\n");
    } else {
        printf("Majeur\n");
    }
    return 0;
}
\end{ccode}

\corrige{de l'exercice 11}
\begin{ccode}
#include <stdio.h>
int main() {
    for (int i = 1; i <= 10; i++) {
        printf("%d ", i);
    }
    printf("\n");
    return 0;
}
\end{ccode}

\corrige{de l'exercice 12}
\begin{ccode}
#include <stdio.h>
int main() {
    for (int i = 2; i <= 20; i += 2) {
        printf("%d ", i);
    }
    printf("\n");
    return 0;
}
\end{ccode}

\corrige{de l'exercice 13}
\begin{ccode}
#include <stdio.h>
int main() {
    int n, somme = 0;
    printf("Entrez un entier n : ");
    scanf("%d", &n);
    for (int i = 1; i <= n; i++) {
        somme += i;
    }
    printf("Somme = %d\n", somme);
    return 0;
}
\end{ccode}

\corrige{de l'exercice 14}
\begin{ccode}
#include <stdio.h>
int main() {
    int n;
    printf("Entrez un entier : ");
    scanf("%d", &n);
    for (int i = 1; i <= 10; i++) {
        printf("%d x %d = %d\n", n, i, n * i);
    }
    return 0;
}
\end{ccode}

\corrige{de l'exercice 15}
La boucle \texttt{for} va de 1 à 5. À chaque itération, on ajoute \texttt{i} à \texttt{s}.
\begin{itemize}
    \item i=1 : s = 0+1 = 1
    \item i=2 : s = 1+2 = 3
    \item i=3 : s = 3+3 = 6
    \item i=4 : s = 6+4 = 10
    \item i=5 : s = 10+5 = 15
\end{itemize}
Sortie : \texttt{s = 15}

\corrige{de l'exercice 16}
\begin{ccode}
#include <stdio.h>
int main() {
    int n, estPremier = 1;
    printf("Entrez un entier : ");
    scanf("%d", &n);
    if (n <= 1) {
        estPremier = 0;
    } else {
        for (int i = 2; i * i <= n; i++) {
            if (n % i == 0) {
                estPremier = 0;
                break;
            }
        }
    }
    if (estPremier) {
        printf("%d est premier.\n", n);
    } else {
        printf("%d n'est pas premier.\n", n);
    }
    return 0;
}
\end{ccode}

\corrige{de l'exercice 17}
\begin{ccode}
#include <stdio.h>
int main() {
    int tab[5] = {10, 20, 30, 40, 50};
    for (int i = 0; i < 5; i++) {
        printf("%d ", tab[i]);
    }
    printf("\n");
    return 0;
}
\end{ccode}

\corrige{de l'exercice 18}
\begin{ccode}
#include <stdio.h>
int main() {
    float notes[10], somme = 0, moyenne;
    printf("Entrez 10 notes : ");
    for (int i = 0; i < 10; i++) {
        scanf("%f", &notes[i]);
        somme += notes[i];
    }
    moyenne = somme / 10;
    printf("Moyenne = %.2f\n", moyenne);
    return 0;
}
\end{ccode}

\corrige{de l'exercice 19}
\begin{ccode}
#include <stdio.h>
int main() {
    int tab[5], max, pos;
    printf("Entrez 5 entiers : ");
    for (int i = 0; i < 5; i++) {
        scanf("%d", &tab[i]);
    }
    max = tab[0];
    pos = 0;
    for (int i = 1; i < 5; i++) {
        if (tab[i] > max) {
            max = tab[i];
            pos = i;
        }
    }
    printf("Le plus grand élément est %d à la position %d.\n", max, pos);
    return 0;
}
\end{ccode}

\corrige{de l'exercice 20}
Compléter la boucle \texttt{for} :
\begin{ccode}
for (i = 0; i < 10; i++) {
    if (tab[i] == v) {
        compteur++;
    }
}
\end{ccode}

\corrige{de l'exercice 21}
\begin{ccode}
int max2(int a, int b) {
    if (a > b) {
        return a;
    } else {
        return b;
    }
}
\end{ccode}

\corrige{de l'exercice 22}
\begin{ccode}
#include <stdio.h>
int estPair(int n) {
    return (n % 2 == 0);
}
int main() {
    int n;
    printf("Entrez un entier : ");
    scanf("%d", &n);
    if (estPair(n)) {
        printf("Pair\n");
    } else {
        printf("Impair\n");
    }
    return 0;
}
\end{ccode}

\corrige{de l'exercice 23}
\begin{ccode}
int sommeTableau(int tab[], int taille) {
    int somme = 0;
    for (int i = 0; i < taille; i++) {
        somme += tab[i];
    }
    return somme;
}
\end{ccode}

\corrige{de l'exercice 24}
\begin{ccode}
#include <stdio.h>
int factoriel(int n) {
    int f = 1;
    for (int i = 1; i <= n; i++) {
        f *= i;
    }
    return f;
}
int main() {
    int n;
    printf("Entrez un entier : ");
    scanf("%d", &n);
    printf("%d! = %d\n", n, factoriel(n));
    return 0;
}
\end{ccode}

\corrige{de l'exercice 25}
\begin{ccode}
#include <stdio.h>
int main() {
    int n;
    printf("Entrez un entier : ");
    scanf("%d", &n);
    printf("Diviseurs de %d : ", n);
    for (int i = 1; i <= n; i++) {
        if (n % i == 0) {
            printf("%d ", i);
        }
    }
    printf("\n");
    return 0;
}
\end{ccode}

\corrige{de l'exercice 26}
\begin{ccode}
#include <stdio.h>
int main() {
    int tab[10], temp;
    printf("Entrez 10 entiers : ");
    for (int i = 0; i < 10; i++) {
        scanf("%d", &tab[i]);
    }
    // Tri à bulles
    for (int i = 0; i < 9; i++) {
        for (int j = 0; j < 9 - i; j++) {
            if (tab[j] > tab[j+1]) {
                temp = tab[j];
                tab[j] = tab[j+1];
                tab[j+1] = temp;
            }
        }
    }
    printf("Tableau trié : ");
    for (int i = 0; i < 10; i++) {
        printf("%d ", tab[i]);
    }
    printf("\n");
    return 0;
}
\end{ccode}

\corrige{de l'exercice 27}
\begin{ccode}
#include <stdio.h>
int main() {
    float montant, remise, aPayer;
    printf("Entrez le montant des achats (FCFA) : ");
    scanf("%f", &montant);
    if (montant > 50000) {
        remise = montant * 0.10;
        aPayer = montant - remise;
        printf("Remise de %.2f FCFA appliquée.\n", remise);
    } else {
        aPayer = montant;
    }
    printf("Montant à payer : %.2f FCFA\n", aPayer);
    return 0;
}
\end{ccode}